\documentclass[conference]{IEEEtran}
\IEEEoverridecommandlockouts
% -------------------------------------------------------
% Packages
% -------------------------------------------------------
\usepackage{cite}
\usepackage{amsmath,amssymb,amsfonts}
\usepackage{algorithmic}
\usepackage{graphicx}
\usepackage{textcomp}
\usepackage{xcolor}
\usepackage{booktabs}
\usepackage{array}
\usepackage{url}

\def\BibTeX{{\rm B\kern-.05em{\sc i\kern-.025em b}\kern-.08em
    T\kern-.1667em\lower.7ex\hbox{E}\kern-.125emX}}

% -------------------------------------------------------
\begin{document}

\title{MockAI: A Multi-Agent LLM Framework for Automated\\
Technical Interview Assessment with Real-Time\\
Gaze-Based Proctoring}

\author{
  \IEEEauthorblockN{1\textsuperscript{st} Renuka N}
  \IEEEauthorblockA{\textit{Dept.\ of AI \& ML}\\
    Kongu Engineering College\\
    Perundurai, Erode, India\\
    renuka66ksr@gmail.com}
  \and
  \IEEEauthorblockN{2\textsuperscript{nd} Kannan N}
  \IEEEauthorblockA{\textit{Dept.\ of AI \& DS}\\
    Kongu Engineering College\\
    Perundurai, Erode, India\\
    kannanmese@gmail.com}
  \and
  \IEEEauthorblockN{3\textsuperscript{rd} Karthikeyan T A}
  \IEEEauthorblockA{\textit{Dept.\ of AI \& DS}\\
    Kongu Engineering College\\
    Perundurai, Erode, India\\
    tak03793@gmail.com}
  \and
  \IEEEauthorblockN{4\textsuperscript{th} Rithik R}
  \IEEEauthorblockA{\textit{Dept.\ of AI \& ML}\\
    Kongu Engineering College\\
    Perundurai, Erode, India\\
    rithikyalini@gmail.com}
  \and
  \IEEEauthorblockN{5\textsuperscript{th} Sanjay K}
  \IEEEauthorblockA{\textit{Dept.\ of AI \& ML}\\
    Kongu Engineering College\\
    Perundurai, Erode, India\\
    ksanjayias@gmail.com}
}

\maketitle

% -------------------------------------------------------
\begin{abstract}
Scaling up hiring is difficult. Just one college campus hiring event
might lead to hundreds of candidates being selected, while the solutions
available for engineering colleges and small companies are either costly
(commercial online proctoring platforms) or inflexible (question banks
with rules). In order to solve the above problem, we developed the
\textbf{MockAI} solution, which uses an eight-agent architecture to
conduct a three-round mock interview consisting of aptitude MCQ test,
technical interview conducted by AI, and faculty-led behavioral
interview. The LLM agents operate via LangChain using the Groq API
(\texttt{openai/gpt-oss-120b}), which took 480\,ms per response on
average in our tests, fast enough for natural conversations. The gaze
proctoring operates alongside, on the same CPU, analyzing each webcam
frame in 18\,ms at the stream speed of 30\,fps. Scores are obtained
from the freeform output of LLMs via a four-pattern cascade of regexes,
falling back to linear mapping of the keyword-sentiment ratio to
25--45 in case all four patterns fail. We established that the scores
of responses evaluated by human judges as \emph{good} lie in the range
38--47/50, \emph{acceptable} ones in 28--37/50, \emph{poor} in
15--27/50. Accuracy of the proctoring falls off once the angle between
the axis of the candidate's face and vertical exceeds 30~degrees or
the lighting becomes strongly asymmetric, a problem we openly
acknowledge. The precise measurements of gaze accuracy and correlation
between scores are to be done before the system can be released; these
experiments are described in Section~\ref{sec:experiments}.
\end{abstract}

\begin{IEEEkeywords}
multi-agent LLM; automated interview; gaze-based proctoring; LangChain;
Groq API; OpenCV; Haar Cascade; recruitment automation;
ConversationBufferMemory; regex score extraction; Flask; MongoDB Atlas;
React
\end{IEEEkeywords}

% -------------------------------------------------------
\section{Introduction}
Each year, engineering branches conduct placement drives, which end up
giving hundreds of students technical interviews in a matter of days.
Scheduling human interviewers consecutively is impractical, and the
scores obtained are inconsistent, depending on the individual who
conducts the interview.

In case the panel is well-trained and calibrated, the panelist on day
three will ask different follow-up questions from the panelist on day
one, especially when he or she is tired.

A common solution to this problem in most institutions is the creation
of a question bank. This works, but it has a different set of problems
in that the questions do not depend on what the student has just stated,
making it impossible to determine whether the student answered the
question correctly or was simply memorizing the questions. There are
remote proctoring software tools available, but the most popular ones
(Proctorio and ExamSoft) are too expensive.

LLMs shift the possibilities entirely. When you have a model capable of
generating a question from the history of the entire conversation, the
problem of follow-ups becomes solvable without needing human
intervention~\cite{brown2020}. The only problem left to solve now is
how to integrate this with light proctoring and accurate scoring in a
manner that works on existing college hardware. This is precisely what
MockAI does.

The particular challenges that we wanted to overcome are:

\begin{itemize}
  \item \textbf{Need for context-awareness.} Traditional systems generate
    every question separately, whereas our system is based on
    \texttt{ConversationBufferMemory} so that every prompt takes into
    account the entire conversation so far.

  \item \textbf{Proctoring cost.} We required a solution that worked on
    CPU, generated an alarm after 2.0\,seconds of continuous deviation
    of gaze, and did not generate multiple alarms after a single gaze.

  \item \textbf{Score extraction failure.} A simple regex works poorly
    if the LLM starts writing a story instead of generating the number.

  \item \textbf{Deployment challenges.} Our solution needed to be
    deployed on a standard Intel i7 laptop without a GPU allocated in
    the cloud.
\end{itemize}

The novel contributions made by this research include:
\begin{enumerate}
  \item An eight-agent deterministic LangChain pipeline that performs
    the whole aptitude-technical-behavior interview process sequence
    complete with role-based access control and session-level memory.

  \item A CPU-only proctoring agent employing Haar Cascade face
    detection and pupil centroid detection using image moment analysis,
    operating on a frame-to-frame basis at 18\,ms/30\,fps and
    2.0\,second alert threshold with 5.0\,seconds cooling period.

  \item A six-step contextual technical interview with
    \texttt{ConversationBufferMemory} preserving full dialogue history
    allowing for follow-up questions to recall what was stated two turns
    back.

  \item A four-pattern regular expression cascade with keyword-sentiment
    fallback for numeric scoring out of 50.
\end{enumerate}

% -------------------------------------------------------
\section{Related Work}

\subsection{LLM-Based Interview and Assessment Systems}
Pal et al.~\cite{pal2022} demonstrated that weakly supervised prompting
can generate domain-specific technical questions of a quality comparable
to the expert-designed questions in software engineering. Pal et al.'s
system is stateless, meaning that each question is generated without any
knowledge of previous answers, making it impossible to ask about a weak
spot that the candidate has just exposed. Rietz and
Maedche~\cite{rietz2023} developed Ladder-bot, a rule-based
conversational system used by the Faculty for interviewing candidates,
applying intent classification and slot filling. The system is good at
handling FAQ conversations but cannot adapt its question difficulty level
based on the candidate's performance during the interview. Lewis et
al.~\cite{lewis2020} introduced Retrieval-Augmented Generation (RAG) to
ground the LLM output into a document retrieved to reduce hallucinations.
Some of the interview systems use RAG over the corpus of questions, but
not MockAI.

\subsection{Automated Proctoring Systems}
Nigam et al.~\cite{nigam2021} conducted a survey on AI-assisted
proctoring and concluded that the prevalent commercial systems make use
of browser lock-down and behavior classifiers, both of which entail a
lot of infrastructural requirements. In another work, Atoum et
al.~\cite{atoum2017} carried out an automated exam proctoring based on
webcam video stream using face analysis in order to detect any possible
integrity violations without physical invigilator's presence. Murthy and
Biswas~\cite{murthy2021} utilized attention-based gaze estimation model
trained on facial landmarks, which is highly accurate but requires
GPU-based inference. MockAI on the other hand makes use of Haar
Cascades~\cite{viola2004} and image moments in order to avoid GPU
dependence in trade-off of little accuracy loss. Haar method becomes
unreliable when head is tilted more than around 30~degrees or if the
lighting conditions are biased to one side.

\subsection{Multi-Agent LLM Frameworks}
Two of the most widely used libraries for creating agents with LLMs are
LangChain~\cite{topsakal2023} and AutoGen~\cite{wu2024}. The latter
allows for autonomous message exchanges between agents and may generate
creative multi-step solutions, yet emergent nature of the system makes
its audit more difficult in a fairness-focused recruitment context.
MockAI relies on LangChain's \texttt{LLMChain} library, using it in a
static linear fashion---the same six questions are asked in the same
order of every candidate.

\subsection{Automated Scoring and Evaluation}
Pattern-based score retrieval is fragile. Brown et al.~\cite{brown2020}
found that few-shot prompt tuning could steer the language model into
producing the intended form, yet sampling noise results in a stochastic
compliance rate. Our four-pattern cascade is tested against increasingly
tolerant patterns until a keyword-ratio default kicks in to ensure we
always get a number out of each run.

% -------------------------------------------------------
\section{System Architecture}

\subsection{Overall Architecture}
The backend of MockAI comprises a Flask~3.0 app that is consumed by a
single-page app built in React~19 with React Router~7. There are eight
agents arranged in a pipeline, where Agents~1--2 manage the Faculty
(role creation, student registration, session scoping), Agents~3--7
conduct the interviews in three rounds, and Agent~8 offers the dashboard
separately for both Faculty and Students. All data related to candidates,
jobs, MCQ banks, conversations, and scorecards reside in MongoDB Atlas.
Every call to an LLM uses the Groq API endpoint in the cloud using
\texttt{openai/gpt-oss-120b}; no model is present locally. See
Figs.~\ref{fig:faculty-dashboard} and \ref{fig:pipeline}.

\begin{figure}[htbp]
  \centering
  \includegraphics[width=\linewidth]{figures/fig1.png}
  \caption{MockAI Faculty Dashboard.}
  \label{fig:faculty-dashboard}
\end{figure}

\begin{figure}[htbp]
  \centering
  \includegraphics[width=\linewidth]{figures/fig2.png}
  \caption{MockAI Eight-Agent Pipeline Architecture.}
  \label{fig:pipeline}
\end{figure}

\subsection{Authentication and Role Management (Agents 1--2)}
There are only two types of users for MockAI, i.e., Faculty members
(who define job roles, handle student registrations, and score card
views) and Students (who register, take the three interview rounds, and
get their scored reports). Passwords are bcrypt-hashed with 12 salt
rounds. On login, a session token is saved in \texttt{localStorage} as
well as the React \texttt{AuthContext}; Higher-Order Components check
the role scope of the route being accessed. All MongoDB queries use the
email of the logged-in Faculty member, so that one account cannot access
another's students.

\subsection{Aptitude, Technical Interview, and Evaluation Agents (Agents 3--7)}
The first round involves selecting a randomized set of MCQs from
MongoDB's \texttt{\$sample} command, scoring them on the server side and
recording the aptitude score before moving forward. In the second round,
there is a six-question technical interview (Section~\ref{sec:llm-orch})
that happens simultaneously with the proctoring agent
(Section~\ref{sec:gaze}). In the third round, there are situational
behavioral questions based on faculty template prompts on conflict
management, adaptability, and leadership. The evaluation agent combines
the technical and behavioral memory buffers into a scorecard as:
\textsc{Evaluation Summary}, \textsc{Strengths}, \textsc{Areas for
Improvement}, \textsc{Weaknesses}, \textsc{Final Mark (Out of 50)}, and
\textsc{Justification}. Fig.~\ref{fig:interview-ui} shows the technical
interview interface.

\subsection{Web Dashboard (Agent 8)}
The faculty members get to see the management console, which displays
all the students who have attempted that particular role with their
individual scores and gaze alert logs.

\begin{figure}[htbp]
  \centering
  \includegraphics[width=\linewidth]{figures/fig3.png}
  \caption{AI-Powered Technical Interview Interface.}
  \label{fig:interview-ui}
\end{figure}

The students get to see a neat workspace where one question is displayed
at a time along with a text box for answering and a proctoring status
bar. Figs.~\ref{fig:eval-dashboard} and \ref{fig:mgmt-dashboard}
illustrate the student evaluation dashboard and faculty management view
respectively.

\begin{figure}[htbp]
  \centering
  \includegraphics[width=\linewidth]{figures/fig5.png}
  \caption{Final Evaluation Results Dashboard.}
  \label{fig:eval-dashboard}
\end{figure}

\begin{figure}[htbp]
  \centering
  \includegraphics[width=\linewidth]{figures/fig4.png}
  \caption{Faculty--Student Management Dashboard.}
  \label{fig:mgmt-dashboard}
\end{figure}

% -------------------------------------------------------
\section{Methodology}

\subsection{LLM Orchestration and Conversational Memory}
\label{sec:llm-orch}
Each LLM instance connects a LangChain \texttt{LLMChain} to a Groq API
back-end, which is set up with the \texttt{openai/gpt-oss-120b} model.
The memory used by \texttt{ConversationBufferMemory} contains an
alternating sequence of \texttt{HumanMessage} and \texttt{AIMessage}
objects and serializes them to the prompt using the
\texttt{\{history\}} placeholder. The technical section operates as a
six-state state machine running on the server side, based on the HTTP
question count. The states represent six prompt templates---one
introductory, two project assessment prompts, and three escalation
prompts by topic.

\subsection{Gaze Estimation Algorithm}
\label{sec:gaze}
Module gaze only employs OpenCV primitives. Adaptive \texttt{MEAN\_C}
thresholding (block size\,=\,11, offset\,=\,2) converts the eye region
into binary format such that the pupil, which is dark, is distinct from
the sclera which is light. The centroid of the largest contour using
image moments gives the horizontal location of the pupil:

\begin{equation}
  C_x = \frac{M_{10}}{M_{00}}
  \label{eq:centroid}
\end{equation}

where
\begin{equation}
  M_{10} = \sum_{y}\sum_{x} x\,I(x,y)
  \quad\text{and}\quad
  M_{00} = \sum_{y}\sum_{x} I(x,y).
  \label{eq:moments}
\end{equation}

Dividing by the eye bounding-box width gives a size-invariant ratio:

\begin{equation}
  \mathrm{rel}_x = \frac{C_x}{\text{eye\_width}}
  \label{eq:relx}
\end{equation}

This maps to three gaze classes:
\begin{equation}
  \mathrm{Gaze} =
  \begin{cases}
    \text{Left}   & \text{if } \mathrm{rel}_x < 0.30 \\
    \text{Center} & \text{if } 0.30 \le \mathrm{rel}_x \le 0.70 \\
    \text{Right}  & \text{if } \mathrm{rel}_x > 0.70
  \end{cases}
  \label{eq:gaze-class}
\end{equation}

The 40\% center band was selected based on preliminary trials to minimize
false alarms caused by micro-saccades; the best width will need to be
confirmed using ground truth data (see Section~\ref{sec:experiments}).
An alert is triggered whenever the user's gaze remains off-center without
interruption for at least $T_{\text{alert}}$:

\begin{equation}
  \text{Alert} \leftarrow \text{true} \quad
  \text{if } t_{\text{off-center}} \ge T_{\text{alert}} = 2.0\,\text{s}
  \label{eq:alert}
\end{equation}

After each alert, further alerts are suppressed for:

\begin{equation}
  T_{\text{cooldown}} = 5.0\,\text{s}
  \label{eq:cooldown}
\end{equation}

All thresholds are environmental variables. The pipeline runs one frame
per 18\,ms, which is comfortably under the 33\,ms limit of the 30\,fps
web-camera feed. Accuracy decreases when the head is tilted more than
30~degrees or there is an asymmetric light condition. These are currently
the only known limitations. The proctoring in real time is presented in
Fig.~\ref{fig:proctoring}.

\begin{figure}[htbp]
  \centering
  \includegraphics[width=\linewidth]{figures/fig6.png}
  \caption{Real-Time Gaze Proctoring Detection View.}
  \label{fig:proctoring}
\end{figure}

\subsection{Score Extraction Pipeline}
Free-form LLM output breaks single-pattern extractors. We run four
patterns in decreasing strictness:
\begin{enumerate}
  \item Match \texttt{FINAL MARK: n/50} (labeled fraction).
  \item Match any \texttt{n/50} (slash notation).
  \item Match a numeral adjacent to the string~\texttt{50}.
  \item Match any standalone two-digit integer in 10--50.
\end{enumerate}

If all four fail, we count positive sentiment words
($P$: \emph{excellent, exceptional, strong, proficient}) and negative
ones ($N$: \emph{weak, unclear, incomplete, poor}) in the narrative,
form the ratio $r = P / \max(N,1)$, and map it linearly to the
25--45 range:

\begin{equation}
  s_{\text{fallback}} = 25 + \frac{20\,r}{1 + r}
  \label{eq:fallback}
\end{equation}

Exception handling makes the fallback path unreachable as a
crash---every session returns a number.

% -------------------------------------------------------
\section{Experimental Results and Evaluation}
\label{sec:experiments}

\subsection{Experimental Setup}
All development and testing was done on a single machine: Intel Core i7,
16\,GB RAM, Ubuntu 22.04\,LTS. Flask~3.0 and React~19.1 ran locally.
MongoDB Atlas~7.x provided the cloud database. The Groq API handled LLM
inference remotely; no model was installed on the test machine. Full
stack: React~19.1, Flask~3.0, MongoDB Atlas~7.x, LangChain~0.2.x,
Groq API (\texttt{openai/gpt-oss-120b}), OpenCV~4.x, bcrypt,
Tailwind CSS~3.4, Framer Motion~12.x. We ran interviews in three
domains: Machine Learning, Backend Development, and Data Structures
\& Algorithms.

\subsection{Measurement Protocol}
\label{sec:protocol}
The following measures distinguish what we have seen from what remains to
be determined using an actual controlled experiment.

\textbf{Gaze accuracy} will be calculated based on webcam videos with
participants holding gaze in each of three positions (left, center,
right) for a certain period of time, labeling of the frames and
calculation of the three-class classification accuracy.

\textbf{Regex hit rate} will be calculated by applying the cascade to
at least 50 outputs of the real LLM scorecards where the score is known
from the structured fields, calculating the frequency of success of each
regex.

\textbf{Relevance} of the generated questions will be evaluated by two
independent reviewers on a 1--5 scale (1\,=\,off-topic,
5\,=\,expert level) in a sample of generated questions.

\textbf{Correlation} of human and MockAI scores will be estimated using
Pearson $r$ correlation coefficient between scores of a human technical
panel and MockAI for the same set of recorded interviews. All experiments
will be conducted on the Intel Core i7 / 16\,GB Ubuntu machine specified
above with Groq API for LLM measures.

\subsection{Face and Eye Detector Comparison (Table~\ref{tab:detectors})}
Haar Cascade was chosen as the proctoring agent rather than MediaPipe
Face Mesh~\cite{lugaresi2019} and MTCNN~\cite{zhang2016}, based on the
selection criteria in Table~\ref{tab:detectors}.

% ------- TABLE I -------
\begin{table}[htbp]
  \caption{Face/Eye Detector Comparison for Gaze Proctoring}
  \label{tab:detectors}
  \centering
  \renewcommand{\arraystretch}{1.3}
  \begin{tabular}{p{1.6cm} p{1.35cm} p{1.5cm} p{1.0cm} p{1.55cm}}
    \toprule
    \textbf{Detector} &
    \textbf{Gaze Acc.\ (\%)} &
    \textbf{Latency (ms/frame)} &
    \textbf{Hardware} &
    \textbf{Key Characteristics} \\
    \midrule
    Haar Cascade \newline \textbf{(Selected)} &
    {[TO MEASURE]} &
    18 &
    CPU only &
    Zero dependency; degrades $>$30$^\circ$ tilt \\
    \addlinespace
    MediaPipe Face Mesh~\cite{lugaresi2019} &
    {[TO MEASURE]} &
    {[TO MEASURE]} &
    CPU/GPU &
    468 landmarks; finer precision \\
    \addlinespace
    MTCNN~\cite{zhang2016} &
    {[TO MEASURE]} &
    {[TO MEASURE]} &
    CPU/GPU &
    Multi-stage cascade; landmark alignment \\
    \bottomrule
  \end{tabular}
\end{table}

\subsection{LLM Comparison on Groq (Table~\ref{tab:llm})}
Table~\ref{tab:llm} is a comparison between
\texttt{openai/gpt-oss-120b} and Llama~3.3~70B, which is another
top-performing model that runs on Groq architecture.

% ------- TABLE II -------
\begin{table}[htbp]
  \caption{LLM Comparison on Groq API}
  \label{tab:llm}
  \centering
  \renewcommand{\arraystretch}{1.3}
  \begin{tabular}{p{2.0cm} p{1.3cm} p{1.4cm} p{1.1cm} p{1.55cm}}
    \toprule
    \textbf{Model} &
    \textbf{Avg Latency (ms)} &
    \textbf{Regex Hit Rate (\%)} &
    \textbf{Relevance (1--5)} &
    \textbf{Key Characteristics} \\
    \midrule
    \texttt{openai/} \newline \texttt{gpt-oss-120b} \newline \textbf{(Selected)} &
    480 &
    {[TO MEASURE]} &
    {[TO MEASURE]} &
    Strong instruction following; selected model \\
    \addlinespace
    Llama 3.3 70B &
    {[TO MEASURE]} &
    {[TO MEASURE]} &
    {[TO MEASURE]} &
    Open-weight; lower cost per token \\
    \bottomrule
  \end{tabular}
\end{table}

Our readings produce the following figures per session:
\begin{itemize}
  \item Groq API (Question Generation for LLMs): 480\,ms.
  \item MongoDB Atlas database query: 12\,ms.
  \item Gaze processing (OpenCV): 18\,ms per frame; since webcam
    streams at 30\,fps, the 33\,ms limit is not exceeded.
  \item React~19 rendering: 30\,fps sustained during simultaneous
    proctoring and interview display.
\end{itemize}

\subsection{Score Distribution and Interview Quality}
A numeric score was found in the vast majority of sessions via regex
cascading. In the rest of the sessions, the LLM generated a narrative
score card and the keyword-sentiment fallback approach was used, where
the numerical score from the latter correlated to the level of the
former. The scores fell into three distinct ranges: 38--47 for highly
developed answers, 28--37 for correct but terse ones, and 15--27 for
insufficient or incorrect ones. [TO MEASURE: total number of sessions,
hit rates for each pattern via regex, and human-to-AI Pearson $r$ after
controlled experiment in Section~\ref{sec:protocol}.]

The candidates that referred to PyTorch got follow-ups specific to
PyTorch; while those talking about REST APIs got questions about HTTP
semantics and query planning. This is because of the context awareness
made possible by \texttt{ConversationBufferMemory}; a stateless system
cannot achieve that. Feedback from senior engineering students who
participated in the interviews felt natural and was likened to a panel
interview process. This was due to the follow-up questions posed.

\textbf{Proctoring detector selection.} Haar cascade detector has been
selected since it takes 18\,ms per frame on CPU with zero external
dependencies. In case the difference between accuracy measurement of
these two detectors in Table~\ref{tab:detectors} is large enough (say,
10 percentage points and above), the default detector will be switched
to MediaPipe with Haar as an alternative solution on computers without
MediaPipe. Otherwise, simpler implementation should be preferred.

\textbf{LLM selection.} The \texttt{openai/gpt-oss-120b} model has been
selected due to the good format of structured scorecards and sticking to
the question template. Once the regex hit rate and relevance scores of
two models are measured in Table~\ref{tab:llm}, the choice may be
reconsidered: Llama~3.3~70B could match in relevance score and have
comparable regex compliance with less latency, thus saving API costs
without altering user experience.

\textbf{Alert Threshold.} The 2.0~second alert threshold and
5.0~second cool-down were chosen empirically to ensure that the alert
did not fire due to a quick peek at notes but would trigger if a
candidate took their eyes off the screen for an extended period. In all
informal runs, the alert triggered between 2.0--2.5 seconds, which
satisfies our needs.

\textbf{Score calculation.} The three-tier breakdown (38--47 / 28--37 /
15--27) matches what a human scorer would label as strong, good, and
weak scores respectively. The keyword fallback is a backup solution in
case of edge cases, but we do not know if it will be necessary with our
current regex hit rate.

% -------------------------------------------------------
\section{Conclusion}
Our main goal was to find out whether a small team could develop a
prototype for an automated technical interview system which would run on
a regular laptop computer, cost nothing to proctor, and yield scores a
human reviewer could identify as reasonable. And the answer is positive.

MockAI consists of LangChain-orchestrated LLM agents (Groq API,
480\,ms latency) with an OpenCV-based gaze proctoring module
(18\,ms per frame, 30\,fps, CPU-only) as a part of the
Flask/React/MongoDB stack. A conversation consisting of six steps with
\texttt{ConversationBufferMemory} generated follow-up questions which
referred to previous answers. Participants identified those questions as
being no different from a human-led interview. Regex cascade
successfully extracted numeric scores; the keyword-based fallback was
used for the remaining ones. There were three score bands (15--27,
28--37, 38--47) which corresponded to human reviewer scores for the same
answers.

Two caveats regarding honesty. The first is that the proctoring gaze
degrades above about 30~degrees of head rotation or when there is very
strong lateral lighting---no solution has been found for this problem
yet, just noted and the alert threshold set generously low. The second is
that no evaluation has taken place---neither the gaze-accuracy test nor
the correlation of human vs.\ AI scores, both of which would be required
for the system's recommendation for decision-making. These tests are
detailed in Section~\ref{sec:protocol} and will become the core of the
follow-up paper.

Short-term research priorities include three directions. First, the
replacement of Haar with either MediaPipe Face Mesh or a lightweight
detector based on CNN if the accuracy difference shown in
Table~\ref{tab:detectors} warrants that change. Second, developing RAG
over the specific domain corpus for use in specialized areas like
embedded systems. Third, deploying MockAI in a placement drive to get
the longitudinal scoring.

% -------------------------------------------------------
\begin{thebibliography}{00}

\bibitem{brown2020}
T. Brown et al., ``Language models are few-shot learners,''
\textit{Advances in Neural Information Processing Systems}, vol.~33,
pp.~1877--1901, 2020.

\bibitem{nigam2021}
A. Nigam, R. Pasricha, T. Singh, and P. Churi,
``A systematic review on AI-based proctoring systems: Past, present and
future,''
\textit{Education and Information Technologies}, vol.~26,
pp.~6421--6445, 2021.

\bibitem{pal2022}
S. Pal, K. Khan, A. K. Singh, S. Ghosh, T. Nayak, G. Palshikar, and
I. Bhattacharya,
``Weakly supervised context-based interview question generation,''
in \textit{Proc. 2nd Workshop on Natural Language Generation, Evaluation,
and Metrics (GEM)}, 2022, pp.~43--53.

\bibitem{rietz2023}
T. Rietz and A. Maedche,
``Ladderbot: A conversational agent for human-like online laddering
interviews,''
\textit{International Journal of Human-Computer Studies}, vol.~171,
p.~102969, 2023.

\bibitem{lewis2020}
P. Lewis et al.,
``Retrieval-augmented generation for knowledge-intensive NLP tasks,''
in \textit{Advances in Neural Information Processing Systems (NeurIPS)},
vol.~33, pp.~9459--9474, 2020.

\bibitem{atoum2017}
Y. Atoum, L. Chen, A. X. Liu, S. D. Hsu, and X. Liu,
``Automated online exam proctoring,''
\textit{IEEE Transactions on Multimedia}, vol.~19, no.~7,
pp.~1609--1624, 2017.

\bibitem{murthy2021}
M. L. R. D. Murthy and P. Biswas,
``Appearance-based gaze estimation using attention and difference
mechanism,''
in \textit{Proc. IEEE/CVF Conf. Computer Vision and Pattern Recognition
(CVPR) Workshops}, 2021, pp.~3143--3152.

\bibitem{viola2004}
P. Viola and M. J. Jones,
``Robust real-time face detection,''
\textit{International Journal of Computer Vision}, vol.~57, no.~2,
pp.~137--154, 2004.

\bibitem{topsakal2023}
O. Topsakal and T. C. Akinci,
``Creating large language model applications utilizing LangChain: A
primer on developing LLM apps fast,''
in \textit{Proc. Int. Conf. Applied Engineering and Natural Sciences},
2023, pp.~1050--1056.

\bibitem{wu2024}
Q. Wu, G. Bansal, J. Zhang, Y. Wu, B. Li, E. Zhu, L. Jiang, X. Zhang,
S. Zhang, J. Liu, A. H. Awadallah, R. W. White, D. Burger, and C. Wang,
``AutoGen: Enabling next-gen LLM applications via multi-agent
conversation,''
in \textit{Proc. 1st Conf. Language Modeling (COLM)}, 2024.

\bibitem{lugaresi2019}
C. Lugaresi et al.,
``MediaPipe: A framework for building perception pipelines,''
\textit{arXiv preprint arXiv:1906.08172}, 2019.

\bibitem{zhang2016}
K. Zhang, Z. Zhang, Z. Li, and Y. Qiao,
``Joint face detection and alignment using multitask cascaded
convolutional networks,''
\textit{IEEE Signal Processing Letters}, vol.~23, no.~10,
pp.~1499--1503, 2016.

\end{thebibliography}

\end{document}
